\mask{55.903}{83.4}{36.805}{12.52}
\mask{94.108}{83.4}{50.1961}{12.52}
\mask{145.704}{83.4}{410.493}{12.52}
\mask{145.704}{97.32}{410.493}{12.52}
\mask{145.704}{111.24}{410.493}{11.92}
\mask{145.704}{124.56}{410.493}{46.6}
\mask{145.704}{172.56}{410.493}{57.88}
\mask{145.704}{231.84}{410.493}{23.59}
\mask{145.704}{256.83}{410.493}{23.44}
\mask{145.704}{281.67}{410.493}{81.04}
\mask{145.704}{364.11}{410.493}{23.56}
\mask{145.704}{389.07}{410.493}{11.92}
\mask{145.704}{402.39}{410.493}{34.98}
\mask{145.704}{438.77}{410.493}{23.56}
\mask{145.704}{463.73}{410.493}{12.04}
\mask{145.704}{477.17}{410.493}{23.44}
\mask{145.704}{502.01}{410.493}{12.04}
\mask{145.704}{515.45}{410.493}{103.99}
\mask{145.704}{620.84}{410.493}{23.56}
\mask{145.704}{645.8}{410.493}{23.44}
\mask{145.704}{670.64}{410.493}{23.436}
\mask{145.704}{695.476}{410.493}{12.04}
\zhbox{54}{64}{504}{17}{14}{\textbf{引脚功能}}
\attext{57}{85}{34}{9}{\centering\textbf{引脚}}
\attext{99}{85}{41}{9}{\centering\textbf{名称}}
\attext{148}{85}{406}{9}{\textbf{说明}}
\attext{148}{99}{406}{9}{\textbf{I\textsuperscript{2}C 数据信号。}}
\attext{148}{112.5}{406}{9}{\textbf{I\textsuperscript{2}C 时钟信号。}}
\attext{148}{126}{406}{9}{\textbf{过流（OC）、过温（OT）、过压（OV）及 PNG 故障事件中断。}默认配置下发生 PNG 事件时，INT 引脚被屏蔽。INT 为开漏输出；发生中断事件时拉低，故障消除后恢复初始状态。IC 禁用时，INT 为开漏状态。}
\attext{148}{174}{406}{9}{\textbf{电流调节脉宽调制（PWM）信号输入。}通过 I\textsuperscript{2}C 寄存器配置输出电流限值（$I_{OUT\_LIMIT}$）后，IPWM 输入信号占空比控制平均输出电流限值。内部经 1 M$\Omega$ 电阻将 IPWM 上拉至 AVDD，以设置 100\% 占空比高电平。建议为 IPWM 另接 100 k$\Omega$ 外部上拉电阻。}
\attext{148}{234}{406}{9}{\textbf{开关电流检测正输入。}连接至电流检测电阻的高侧（HS）。}
\attext{148}{259}{406}{9}{\textbf{开关电流检测负输入。}连接至电流检测电阻的低侧（LS）。}
\attext{148}{283.5}{406}{9}{\textbf{开关频率（$f_{SW}$）位默认值设置。}启动期间，FS 引脚的两种电平设置默认频率：第一种为 $0.51\times V_{AVDD}$ 至 $0.68\times V_{AVDD}$，FSW 位设为 00（默认 200 kHz）；第二种至少为 $0.74\times V_{AVDD}$，FSW 位设为 10（默认 400 kHz）。FS 不得悬空或接地。各电平设置见第 9 页“电气特性”。启动后改变 FS 电平不会改变 FSW 位；启动后可通过 I\textsuperscript{2}C 接口设置 FSW 位。}
\attext{148}{366}{406}{9}{\textbf{内部控制电路偏置电源。}使用不小于 2.2 $\mu$F 的电容对 AVDD 去耦。}
\attext{148}{391}{406}{9}{\textbf{模拟地。}}
\attext{148}{404.5}{406}{9}{\textbf{电流监测输出。}IMON 输出 IAVGP 与 IAVGN 引脚间信号所对应的电压信号。MP2984 通过电流检测电阻检测平均负载电流，并输出该信号。}
\attext{148}{441}{406}{9}{\textbf{内部误差放大器（EA）输出。}在 COMP 与 AGND 之间串联电容和电阻，以进行环路补偿。}
\attext{148}{465.5}{406}{8.8}{\textbf{输出电压（$V_{OUT}$）反馈。}在 VOUT 与 FB 之间连接电阻分压器。}
\attext{148}{479}{406}{9}{\textbf{软启动（SS）设置。}在 SS 引脚接外部电容；SS 还设置打嗝模式关断时间。}
\attext{148}{504}{406}{9}{\textbf{不连接。}}
\attext{148}{517.5}{406}{9}{\textbf{ILIM 位默认值设置。}启动期间，ILIM 引脚的两种电平设置默认电流限值（$I_{LIMIT}$）。第一种电平使 ILIM 位为 001（默认 32 mV），第二种为 011（默认 45 mV），第三种为 101（默认 56 mV），第四种为 111（默认 68 mV）。默认配置下悬空 ILIM，将 ILIMT 位设为 111。各电平设置见第 9 页“电气特性”。启动后改变 ILIM 电平不会改变 ILIM 位；启动后可通过 I\textsuperscript{2}C 接口设置 ILIM 位。}
\attext{148}{623}{406}{9}{\textbf{平均电流限值检测负输入。}将 IAVGN 和 IAVGP 接至输出轨正端，以设置输出电流限值 $I_{OUT\_LIMIT}$。}
\attext{148}{647.5}{406}{9}{\textbf{平均电流限值检测正输入。}将 IAVGN 和 IAVGP 接至输出轨正端，以设置输出电流限值 $I_{OUT\_LIMIT}$。}
\attext{148}{672.5}{406}{9}{\textbf{输出电压检测输入。}VOUT 根据 VCC 电源逻辑向 VCC 供电；将其连接至输出电容（$C_{OUT}$）。}
\attext{148}{697.5}{406}{9}{\textbf{变换器升压开关节点。}连接 SWD 源极与 SWC 漏极。}
